\subsection{Order, Noncommutativity, and Associativity}

\begin{corebox}[Changing the order changes the question]
The product $AB$ combines rows of $A$ with columns of $B$. Reversing the order
produces a different interaction, so matrix multiplication is not generally
commutative.
\end{corebox}

\begin{examplebox}[The same matrices in the opposite order]
For
\[
A=\begin{pmatrix}1&3\\-2&4\end{pmatrix},
\qquad
B=\begin{pmatrix}2&0\\5&-1\end{pmatrix},
\]
we already found
\[
AB=\begin{pmatrix}17&-3\\16&-4\end{pmatrix}.
\]
In the opposite order,
\[
BA=
\begin{pmatrix}2&0\\5&-1\end{pmatrix}
\begin{pmatrix}1&3\\-2&4\end{pmatrix}
=
\begin{pmatrix}2&6\\7&11\end{pmatrix}.
\]
Hence $AB\neq BA$.
\end{examplebox}

\begin{remarkbox}[Commutativity is a property to check]
Some pairs of matrices do commute, but this is never automatic. Whenever the
order matters in an argument, commutativity must be proved from the structure of
the matrices involved.
\end{remarkbox}

\begin{theorembox}[Associativity]
Whenever all products are defined,
\[
(AB)C=A(BC).
\]
Associativity allows products in a fixed order to be grouped differently. It does
not allow the matrices themselves to be rearranged.
\end{theorembox}

\begin{interpretationbox}[Grouping versus reordering]
Parentheses determine which multiplication is carried out first; the written
order determines which transformations or interactions occur first. Associativity
changes grouping only, while commutativity would change order.
\end{interpretationbox}
